\documentclass[10pt,journal,compsoc]{IEEEtran}

\usepackage[utf8]{inputenc}
\usepackage{amsmath,amssymb,amsfonts}
\usepackage{algorithmic}
\usepackage{graphicx}
\usepackage{textcomp}
\usepackage{xcolor}
\usepackage{booktabs}
\usepackage{cite}
\usepackage{url}
\usepackage{hyperref}

\hypersetup{
    colorlinks=true,
    linkcolor=blue,
    citecolor=blue,
    urlcolor=blue
}

\begin{document}

\title{GestureCAD: An Interactive, Vision-Driven 3D CAD Modeling System with Swept Polyhedral Geometry and Multi-Surface Sketching}

\author{Alish~Laimayum
\IEEEcompsocitemizethanks{\IEEEcompsocthanksitem Department of Computer Science \& Engineering.\protect\\
E-mail: alish.laimayum@example.com}
\thanks{Manuscript received October 2026.}}

\markboth{IEEE Transactions on Visualization and Computer Graphics,~Vol.~XX, No.~X, 2026}%
{Laimayum: GestureCAD: Vision-Driven 3D CAD Modeling System}

\IEEEtitleabstractindextext{%
\begin{abstract}
Traditional Computer-Aided Design (CAD) workflows remain constrained to classical WIMP (Windows, Icons, Menus, Pointer) interfaces, requiring specialized 2D hardware peripherals and indirect multi-view projection manipulation to author 3D spatial geometry. While immersive Virtual and Augmented Reality (VR/AR) solutions provide spatial input, they impose substantial economic burdens, cumbersome head-mounted displays, and fatigue-inducing interaction paradigms, while predominantly producing open, non-manifold surface approximations. In this paper, we present \textbf{GestureCAD}, a contactless, web-based 3D CAD modeling system driven by commodity monocular RGB hand tracking. We introduce a decoupled client-edge architecture that bridges an asynchronous edge computer vision pipeline (OpenCV/MediaPipe) with a zero-reconciliation React Three Fiber WebGL rendering engine.

To resolve spatial ambiguity and hand tremor without tactile feedback, GestureCAD incorporates: (1) an anatomical scale-invariant normalized metric space using wrist-to-MCP Euclidean baselines ($0.0\%$ theoretical CV across camera depths under perspective projection); (2) a dual-threshold hysteresis state automaton reducing spurious pinch state flips by an aggregate mean of $85.55\%$ [$95\%$ CI: $84.01\%$–$87.10\%$]; (3) a dynamic velocity-adaptive smoothing filter with dead-zone noise gating reducing RMS cursor jitter by a mean of $52.43\%$ [$95\%$ CI: $52.16\%$–$52.71\%$]; (4) a 3D geometric snapping and heuristic shape-regularization engine capturing $84.48\%$ of targeting queries and reducing positional error by $79.74\%$ [$95\%$ CI: $77.28\%$–$82.20\%$]; and (5) direct-manipulation linear extrusion and rotational sweeping engines that synthesize closed polyhedral solid meshes. Furthermore, a surface-attached coordinate mapping pipeline constructs local orthonormal frames $\mathbf{R} \in SO(3)$ on existing 3D boundaries, permitting hierarchical sketch-on-face operations directly in 3D space. Comprehensive empirical evaluations demonstrate software CAD execution within $0.1994\text{ ms}$ mean CPU time per frame on host hardware, 100\% compliance with 2-manifold topological invariants ($\chi = V - E + F = 2$, zero boundary edges, zero non-manifold edges) on genus-0 solids, exact Divergence Theorem volumetric integration ($<0.001\%$ numerical discrepancy), and native export to Wavefront OBJ, AutoCAD DXF, and stereolithography STL.
\end{abstract}

\begin{IEEEkeywords}
Computer-Aided Design, Gesture-Based Interaction, Monocular Hand Tracking, Polyhedral Solid Modeling, Topological Invariants, Divergence Theorem, WebGL, Zero-Reconciliation Rendering.
\end{IEEEkeywords}}

\maketitle
\IEEEdisplaynontitleabstractindextext
\IEEEpeerreviewmaketitle

\section{Introduction}
\label{sec:introduction}

\IEEEPARstart{S}{ince} Sutherland's seminal introduction of \textit{Sketchpad} in 1963~\cite{sutherland1963sketchpad}, Computer-Aided Design (CAD) systems have become the foundational infrastructure for mechanical engineering, aerospace fabrication, architecture, and additive manufacturing. Over the subsequent six decades, geometric modeling engines have matured from rudimentary 2D wireframe drafters to mathematically rigorous Boundary Representation (B-Rep) and Constructive Solid Geometry (CSG) kernels~\cite{mantyla1988,hoffmann1989geometric}. 

However, the foundational interaction paradigm for desktop CAD has remained fundamentally anchored to the classical WIMP (Windows, Icons, Menus, Pointer) interface: three-dimensional spatial constructs are continuously translated through two-dimensional planar input peripherals (mouse, trackball, graphics tablet) and visually multiplexed across disjoint orthogonal 2D projection viewports (Top, Front, Right, and Isometric). While experienced CAD practitioners develop remarkable cognitive fluency with this projection mapping, the indirectness introduces non-trivial cognitive load for novice creators, design ideation, and rapid spatial prototyping.

To bridge this spatial rift, the human-computer interaction (HCI) and computer graphics communities have extensively explored immersive spatial authoring paradigms, particularly using stereoscopic Virtual Reality (VR) and Augmented Reality (AR) headsets paired with six-degree-of-freedom ($6$-DoF) optical spatial controllers~\cite{gravitysketch,tiltbrush}. Although immersive spatial modeling systems offer intuitive spatial input, they impose distinct practical trade-offs:
\begin{enumerate}
    \item \textbf{Setup and Peripheral Friction}: Commercial head-mounted displays (HMDs) require significant financial investments, room-scale tracking calibration, and physical tethering or dedicated compute hardware.
    \item \textbf{Physiological Fatigue}: Continuous mid-air arm suspension in unconstrained spatial environments rapidly induces acute musculoskeletal exhaustion, extensively documented in spatial interaction literature as ``Gorilla Arm'' syndrome~\cite{song2014air}.
    \item \textbf{Geometric and Topological Inadequacy}: Existing spatial sketching tools predominantly generate open, non-manifold polygonal ribbons, artistic triangle soups, or loose spline bundles. While visually evocative for conceptual styling, these representations lack topological closure, surface orientation consistency, and volumetric definiteness, rendering them unusable for downstream engineering tasks such as finite element analysis (FEA), stereolithography 3D printing (STL), or CNC toolpath generation.
\end{enumerate}

Concurrently, the emergence of deep learning-based optical hand-pose estimation frameworks---notably Google MediaPipe Hands~\cite{zhang2020mediapipe,lugaresi2019mediapipe} and TensorFlow.js HandPose~\cite{ACP2020handpose}---has democratized real-time articulated 3D hand tracking on commodity monocular RGB webcams, requiring zero specialized infrared sensors (e.g., Leap Motion~\cite{leapmotion2013}) or wearable markers. 

Nevertheless, adapting raw landmark coordinate streams to precision CAD authoring presents formidable signal-processing and systems-engineering hurdles:
\begin{itemize}
    \item \textbf{High-Frequency Physiological Tremor}: Inherent involuntary neuromuscular micro-tremor ($\sigma \approx 1$--$5\text{ mm}$) and optical landmark localization noise produce severe spatial jitter, preventing straight-line authoring and accurate vertex closure.
    \item \textbf{Distance-Dependent Perspective Scaling}: In monocular pinhole optics, apparent hand pixel dimensions vary inversely with camera distance ($Z_{\text{cam}}$), causing fixed pixel-distance gesture thresholds to fail as users naturally adjust their seating position.
    \item \textbf{Limit-Cycle Decision Boundary Chatter}: When a user's hand hovers near the activation threshold of a discrete gesture (e.g., pinch-to-draw), signal noise triggers rapid, uncontrolled state flips between drawing and idle states.
    \item \textbf{Client-Side Web Performance Bottlenecks}: Modern web applications built on declarative UI component trees (e.g., React) incur garbage-collection pauses and Virtual DOM reconciliation overhead when processing high-frequency spatial tracking streams ($30$--$60\text{ Hz}$).
\end{itemize}

\begin{figure}[t]
    \centering
    \fbox{%
        \parbox[c][5.2cm][c]{0.92\linewidth}{%
            \centering\small\sffamily
            \textbf{Figure 1: Application Overview Placeholder}\\
            \vspace{0.3cm}
            \textit{Intended Screenshot:}\\
            GestureCAD user interface displaying the interactive 3D WebGL viewport,\\
            active working plane grid, real-time hand-tracking video overlay with skeletal joints,\\
            continuous drawing cursor, gesture classification state banner, and CAD tools.\\
            \vspace{0.3cm}
            \footnotesize Expected file: \texttt{paper/figures/fig1_overview.png}
        }%
    }
    \caption{The GestureCAD interactive interface. The system captures monocular webcam input, performs kinematic conditioning to filter hand tremor, and renders topologically validated solid geometry in an interactive WebGL viewport.}
    \label{fig:overview}
\end{figure}

To address these challenges, we present \textbf{GestureCAD}, a contactless, web-based 3D CAD modeling application and swept-geometry engine driven entirely by monocular hand gestures. GestureCAD couples a Python-based edge vision microservice with a Three.js/WebGL client-side rendering architecture, transforming natural hand gestures into precision geometric construction commands on standard desktop hardware without proprietary depth sensors.

\subsection{Research Objectives and Scope}
This paper does not claim to introduce a new mathematical solid modeling paradigm or a novel machine-learning architecture for hand pose estimation. Rather, this work addresses a practical systems-engineering problem: \textit{How can commodity monocular hand tracking be effectively conditioned, stabilized, and coupled to an in-browser solid geometry pipeline to achieve interactive, topologically verified 3D modeling without specialized hardware?}

The specific engineering scope encompasses:
\begin{enumerate}
    \item Formulating an integrated kinematic conditioning pipeline tailored to spatial CAD input, including anatomical scale normalization, dual-threshold hysteresis debouncing, dead-zone gating, and dynamic tremor smoothing.
    \item Designing a decoupled client-edge architecture that streams tracking telemetry over WebSockets and updates WebGL vertex buffers directly in memory, bypassing Virtual DOM reconciliation.
    \item Implementing direct-manipulation linear extrusion and rotational sweeping engines that generate polyhedral meshes with automated topological invariant checks ($\chi = 2$) and Divergence Theorem volume calculations.
    \item Establishing a surface-attached local coordinate mapping pipeline ($SO(3)$ orthonormal frames) to support multi-surface sketch-on-face interactions.
    \item Conducting reproducible empirical benchmarks and parameter sweeps evaluating latency, noise attenuation, state stability, and geometric validity.
\end{enumerate}

\subsection{Summary of Contributions}
The specific contributions of this work are:
\begin{itemize}
    \item \textbf{An Integrated Kinematic Conditioning Pipeline}: An anthropometrically normalized pinch metric invariant to camera distance under perspective projection ($0.0\%$ CV pinhole formulation), combined with dual-threshold hysteresis reducing boundary state chatter by a mean of $85.55\%$ [$95\%$ CI: $84.01\%$–$87.10\%$] across 180 swept trials, and dynamic smoothing reducing RMS cursor tremor by $52.43\%$ [$95\%$ CI: $52.16\%$–$52.71\%$].
    \item \textbf{A Decoupled Client-Edge WebCAD System}: An architecture that bridges a multi-threaded edge computer vision capture loop with a zero-reconciliation Three.js rendering engine, executing software CAD operations in $0.1994\text{ ms}$ mean CPU time per frame on host hardware.
    \item \textbf{Swept Polyhedral Geometry with Topological Invariant Verification}: A custom mesh solid generator supporting linear extrusion, rotational sweeping (with Rodrigues' rotation and on-axis vertex welding), exact Divergence Theorem volumetric integration, and automated verification of closed 2-manifold invariants ($\chi = 2$) for genus-0 solids, with explicit characterization of cap triangulation constraints on concave profiles.
    \item \textbf{Multi-Surface Coordinate Mapping}: An orthonormal frame derivation pipeline $\mathbf{R} \in SO(3)$ enabling dynamic sketch-plane attachment on planar polyhedron faces and cylindrical boundaries.
    \item \textbf{A Reproducible Benchmarking Suite}: An open-source empirical evaluation suite classifying benchmarks into physical real-system measurements, stochastic parameter sweeps with 95\% confidence intervals, and analytical derivations.
\end{itemize}

The remainder of this paper is organized as follows: Section~\ref{sec:related_work} situates GestureCAD within prior literature on vision-based interaction, spatial CAD, and web geometry engines. Section~\ref{sec:architecture} describes the system architecture and communication pipeline. Section~\ref{sec:methodology} presents the mathematical formulation of kinematic conditioning and coordinate projection. Section~\ref{sec:geometry} details geometric construction, validation, and face coordinate frames. Section~\ref{sec:experiments} reports experimental results and parameter sweeps. Section~\ref{sec:discussion} analyzes engineering trade-offs and threats to validity. Section~\ref{sec:limitations} explicitly documents system limitations, and Section~\ref{sec:conclusion} concludes with future research directions.

\section{Related Work}
\label{sec:related_work}

The design of GestureCAD intersects three principal domains of research: vision-based hand tracking and signal conditioning, spatial and immersive CAD interaction paradigms, and browser-hosted computational geometry engines.

\subsection{Vision-Based Hand Tracking and Kinematic Conditioning}
Optical markerless hand pose estimation has evolved from early color-segmentation and depth-map skeletal tracking algorithms to deep learning-based regression architectures. Early spatial interaction systems, such as FAAST~\cite{suma2011faast}, relied on structured-light infrared depth sensors (e.g., Microsoft Kinect) to infer coarse arm and upper-body skeletons, but lacked the spatial resolution required for fine multi-finger articulation. Ultraleap's Leap Motion controller~\cite{leapmotion2013} achieved sub-millimeter finger tracking using dedicated desk-mounted stereo infrared cameras and LED illumination; however, its adoption remains constrained by the requirement for specialized external hardware and a restricted optical tracking volume.

The introduction of Google MediaPipe Hands~\cite{zhang2020mediapipe,lugaresi2019mediapipe} marked a critical transition, enabling monocular 2.5D landmark regression directly from commodity RGB video streams at real-time frame rates ($> 25\text{ FPS}$). MediaPipe utilizes a two-stage pipeline: a single-shot detector identifies hand palm boundaries, followed by a regression network estimating the 3D coordinates of 21 articulated joint landmarks. Grishchenko et al.~\cite{ACP2020handpose} subsequently demonstrated in-browser execution via TensorFlow.js.

However, commodity vision streams provide raw coordinate estimates with notable physiological and sensor noise. Unfiltered hand coordinates exhibit micro-tremor ($\sigma \approx 1$--$5\text{ mm}$), optical quantization jitter, and perspective foreshortening. While adaptive filtering algorithms such as the $1\text{\euro}$ Filter~\cite{casiez2012oneeuro} and double exponential smoothing~\cite{laviola2003double} have successfully minimized jitter in general 2D pointing tasks, spatial CAD sketching imposes unique demands: spatial strokes require dead-band noise gating during resting periods, velocity-adaptive smoothing during active contour authoring, distance-invariant gesture thresholds, and limit-cycle hysteresis to prevent pen-state chatter. In GestureCAD, we integrate these requirements into a specialized kinematic conditioning pipeline.

\subsection{Spatial, Gesture-Driven, and Immersive 3D CAD}
The aspiration to interact with CAD models through direct spatial gestures extends back several decades. An important milestone in vision-driven geometric modeling was \textit{Shape-It-Up} by Vinayak et al.~\cite{vinayak2013shapeitup}, which utilized depth cameras to interpret hand gestures for interactive bending, scaling, and twisting of generalized cylindrical shapes. While demonstrating the intuitive appeal of spatial manipulation, the system was restricted to continuous deformation of pre-existing cylindrical topologies and did not support arbitrary 2D sketch-to-solid extrusion, face-local sketching, or closed polyhedral boundary representations.

In the realm of immersive computing, commercial Virtual Reality tools such as Gravity Sketch~\cite{gravitysketch} and Google Tilt Brush~\cite{tiltbrush} introduced sweeping six-degree-of-freedom brush strokes in 3D space. While highly effective for conceptual automotive styling and artistic ideation, these immersive authoring tools fundamentally trade engineering rigor for artistic fluidity: they primarily construct open, non-manifold polygonal ribbons, non-orientable surface meshes, or loose spline networks. Such representations do not satisfy the mathematical criteria of solid models (orientability, closure, and 2-manifold topology), preventing direct downstream volumetric calculation, boolean evaluation, or physical fabrication via standard slicing software.

Conversely, established desktop CAD systems---such as Dassault Syst\`emes SolidWorks~\cite{solidworks}, Autodesk AutoCAD~\cite{autocad}, and the open-source FreeCAD platform~\cite{freecad}---enforce rigorous Boundary Representation (B-Rep) and Constructive Solid Geometry (CSG) standards~\cite{mantyla1988,hoffmann1989geometric}. These platforms maintain complete topological data structures (e.g., winged-edge, radial-edge, or half-edge meshes), exact NURBS surface definitions, parametric feature dependency trees, and bidirectional geometric constraint solvers. However, their interaction paradigm remains strictly anchored to the 2D mouse and keyboard. Novice users must mentally decompose spatial features into disjoint planar projections, introducing a cognitive barrier during initial concept sketching. 

GestureCAD explores an intermediate design space: coupling the spatial directness of natural hand gestures with explicit topological validation to produce closed, watertight polyhedral solids directly from planar sketch profiles, without requiring expensive VR peripherals.

\subsection{Web-Based Computational Geometry and Real-Time Graphics}
With the advent of WebGL and WebAssembly (WASM), geometric modeling has increasingly migrated to browser environments. Three.js~\cite{threejs} provides a high-level scenegraph and GPU shader abstractions, while React Three Fiber (R3F)~\cite{r3f} enables declarative component composition. Onshape demonstrated the commercial viability of cloud-native CAD, executing server-side geometric kernel computations with client-side WebGL streaming.

For client-side geometry evaluation, recent efforts have focused on porting C++ CAD engines to WebAssembly: OpenCASCADE.js~\cite{opencascadejs} compiles the Open CASCADE B-Rep kernel to WASM, enabling full NURBS evaluation and STEP/IGES import in JavaScript. Similarly, ManifoldCAD~\cite{manifoldcad} provides high-performance, guaranteed-manifold boolean operations on triangle meshes via spatial decomposition. 

However, integrating spatial tracking streams into declarative web frameworks presents performance bottlenecks: streaming 30--60 Hz landmark coordinates through React's Virtual DOM reconciler incurs significant garbage collection and reconciliation latency. GestureCAD addresses this bottleneck through a dual-frequency architecture that routes high-frequency position data directly to mutable GPU typed arrays via Three.js \texttt{useFrame} hooks, while reserving reactive state updates for transactional CAD operations.

\begin{table*}[t]
\caption{Systematic comparison of GestureCAD against representative CAD systems, immersive spatial authoring tools, and gesture-modeling prototypes across functional and geometric dimensions.}
\label{tab:related_work}
\centering
\scriptsize
\begin{tabular}{p{2.5cm}p{1.8cm}p{1.8cm}p{1.8cm}p{2.2cm}p{2.2cm}p{2.0cm}p{1.8cm}}
\toprule
\textbf{Dimension / Feature} & \textbf{Commercial CAD} (e.g., SolidWorks~\cite{solidworks}) & \textbf{Open-Source CAD} (e.g., FreeCAD~\cite{freecad}) & \textbf{Immersive VR CAD} (e.g., Gravity Sketch~\cite{gravitysketch}) & \textbf{Gesture Prototype} (e.g., Shape-It-Up~\cite{vinayak2013shapeitup}) & \textbf{Web Frameworks} (e.g., OpenCASCADE.js~\cite{opencascadejs}) & \textbf{Web Cloud CAD} (e.g., Onshape) & \textbf{GestureCAD} (This Work) \\
\midrule
\textbf{Primary Input Modality} & 2D Mouse + Keyboard & 2D Mouse + Keyboard & 6-DoF Tracked VR Controllers & Depth Camera Hand Gestures & Mouse / Scripting APIs & Mouse + Keyboard & Monocular RGB Hand Gestures \\
\textbf{Hardware Peripherals} & Standard Desktop PC & Standard Desktop PC & Dedicated VR HMD + Controllers & IR Depth Sensor (Kinect) & Standard Web Browser & Standard Web Browser & Commodity USB Webcam \\
\textbf{Deployment Model} & Native Desktop ($>10\text{ GB}$) & Native Desktop & Native VR App (SteamVR/Quest) & Standalone Lab Prototype & In-Browser WASM Library & Browser Client + Cloud Kernel & Client-Edge (Python + WebGL) \\
\textbf{Geometric Representation} & Exact NURBS B-Rep (Parasolid) & Exact NURBS B-Rep (OCCT) & Sub-D Surfaces / Open Polygon Ribbons & Deformed Mesh Generalized Cylinders & Exact NURBS B-Rep (OCCT) & Exact NURBS B-Rep & Triangulated Polyhedral Surface Mesh \\
\textbf{Topological Validation} & Complete Manifold \& Non-Manifold Engine & OCCT Kernel Validity Checks & None / Not Enforced (Open Boundary Common) & Surface Vertex Constraints & OCCT Topological Validity Checks & Complete Industrial Kernel Validation & Automated Invariant ($\chi = 2$, $\partial \Omega = 0$, Deg. 2) \\
\textbf{Extrusion \& Revolve} & Full Feature Tree (Draft, Thin, Sweep) & Pad \& Revolution Features & Spatial Curve Lofting / Stroke Sweeping & Axial Scaling / Bending of Primitives & Algorithmic Swept Surfaces & Full Parametric Features & Linear Extrusion + Revolve (Rodrigues) \\
\textbf{Sketching on Faces} & Planar Faces \& Datum Planes & Planar Faces & Freehand Spatial Snapping & Not Supported / Not Documented & Algorithmic Basis Alignment & Planar Faces & Local $SO(3)$ Tangent Coordinate Frames \\
\textbf{Tremor Compensation} & N/A (Mechanical Mouse Friction) & N/A (Mechanical Mouse Friction) & Controller Smoothing Filters & Basic Spatial Averaging & N/A & N/A & 3-Stage Pipeline (Deadzone + EMA + Hysteresis) \\
\textbf{CSG Booleans} & Full Solid Booleans & Full Solid Booleans & Mesh Intersections (Sub-D) & Cylindrical Union / Merging & Complete Boolean Operations & Complete Solid Booleans & Not Implemented (Future Work) \\
\textbf{Export Capabilities} & STEP, IGES, Parasolid, STL, DXF & STEP, IGES, STL, OBJ, DXF & OBJ, FBX, IGES (Pro) & STL & STEP, IGES, BREP, STL & STEP, Parasolid, STL, DXF & Wavefront OBJ, AutoCAD DXF, STL, JSON \\
\textbf{Primary Design Focus} & Industrial Production \& Manufacturing & Parametric Open Engineering & Aesthetic Industrial Concept Styling & Interactive Primitive Shape Expression & In-Browser Geometry Computation & Enterprise Collaborative Cloud CAD & Accessible Gesture-Driven Spatial Prototyping \\
\bottomrule
\end{tabular}
\end{table*}

\subsection{Comparative Analysis and System Positioning}
Table~\ref{tab:related_work} synthesizes the architectural, geometric, and functional characteristics of GestureCAD alongside established desktop CAD platforms, immersive spatial modeling environments, gesture research prototypes, and web geometry frameworks.

Rather than claiming algorithmic superiority over mature industrial systems, this comparison clarifies the specific trade-offs and engineering boundaries of GestureCAD:
\begin{enumerate}
    \item \textbf{Geometric Kernel Maturity}: Commercial desktop CAD (SolidWorks) and open-source systems (FreeCAD) provide complete parametric feature trees, exact NURBS curves, fillets, chamfers, and boolean CSG operations backed by multi-decade industrial kernels (Parasolid, Open CASCADE). GestureCAD does not attempt to reproduce this extensive feature breadth. Its geometric core is a specialized polyhedral mesh generator with automated 2-manifold validation ($\chi = 2$), designed specifically for direct-manipulation linear and rotational sweeping.
    \item \textbf{Input Modality and Setup Burden}: Desktop CAD requires seated WIMP interaction, and immersive VR CAD (Gravity Sketch) demands high-cost headsets and induces musculoskeletal fatigue. In contrast, GestureCAD operates on commodity monocular webcams, eliminating hardware barriers while preserving spatial directness.
    \item \textbf{Signal Conditioning}: While VR controllers rely on onboard IMUs and optical tracking constellations for millimeter accuracy, monocular computer vision landmark streams require specialized kinematic filtering. GestureCAD explicitly addresses this with an integrated scale-normalized hysteresis and velocity-adaptive tremor attenuation stack.
    \item \textbf{Deployment Model}: Unlike browser-based cloud CAD (Onshape) which hosts its kernel on enterprise server clusters, GestureCAD distributes execution between a local edge vision daemon and a client-side WebGL engine.
\end{enumerate}
In summary, GestureCAD occupies a distinct and practical niche: an accessible, webcam-driven spatial modeling environment that combines the fluid affordance of hand gestures with mathematically verified polyhedral geometry.

\section{System Architecture}
\label{sec:architecture}

GestureCAD adopts a decoupled client-edge architecture designed to balance real-time computer vision inference throughput with interactive 60~FPS graphics rendering. As illustrated in Figure~\ref{fig:architecture}, the system separates high-frequency optical tracking from the interactive geometric modeling client across a lightweight local WebSocket network layer.

\begin{figure*}[t]
    \centering
    \fbox{%
        \parbox[c][5.8cm][c]{0.92\linewidth}{%
            \centering\small\sffamily
            \textbf{Figure 2: System Architecture Diagram Placeholder}\\
            \vspace{0.3cm}
            \textit{Intended Diagram:}\\
            Detailed system topology showing dataflow across the two primary tiers:\\
            (1) \textbf{Edge Vision Service (Python)}: Commodity Webcam $\to$ OpenCV Capture Thread $\to$ MediaPipe Hands Inference (21 Landmarks) $\to$ Kinematic Signal Conditioning (Anatomical Normalization, Hysteresis FSM, Deadzone Gating, EMA) $\to$ JSON Telemetry Serializer $\to$ FastAPI WebSocket Server.\\
            (2) \textbf{WebSocket Transport}: Full-duplex JSON streaming on \texttt{ws://localhost:8000/ws/tracking}.\\
            (3) \textbf{Browser CAD Engine (React 18 + Three.js)}: Dual-path execution splitting into:\\
            -- \textit{High-Frequency Render Path}: \texttt{useFrame()} direct typed array (\texttt{Float32Array}) buffer mutations bypassing Virtual DOM $\to$ WebGL GPU Viewport;\\
            -- \textit{Transactional CAD State Path}: Ray-Plane Unprojection $\to$ Snapping \& Regularization $\to$ Zustand Store $\to$ Command History (Undo/Redo) $\to$ Geometry Engine (\texttt{ExtrusionEngine}, \texttt{RevolveEngine}, \texttt{SolidMesh}, \texttt{FaceEnumerator}) $\to$ Exporters (OBJ, DXF, STL, JSON).\\
            \vspace{0.3cm}
            \footnotesize Expected file: \texttt{paper/figures/fig2_architecture.png}
        }%
    }
    \caption{The decoupled client-edge architecture of GestureCAD. Real-time vision processing executes in a dedicated Python background process, while the interactive CAD application and WebGL rendering engine execute client-side in the web browser, communicating via local WebSockets.}
    \label{fig:architecture}
\end{figure*}

\subsection{Edge Vision Microservice}
The vision tracking subsystem is implemented as a multi-threaded Python 3.10+ service utilizing OpenCV for video frame acquisition and Google MediaPipe Hands for articulated skeletal inference. The subsystem executes an asynchronous event loop powered by FastAPI and Uvicorn:

\begin{enumerate}
    \item \textbf{Asynchronous Frame Acquisition}: A dedicated capture thread polls the host video device (defaulting to camera index 0) via \texttt{cv2.VideoCapture}. In empirical testing on host hardware, the acquisition loop achieved an average capture rate of $26.68\text{ FPS}$ with a standard $720\text{p}$ ($1280 \times 720$) resolution sensor.
    \item \textbf{Hand Pose Estimation}: Each captured frame is converted from BGR to RGB color space and processed by MediaPipe Hands configured with \texttt{static\_image\_mode=False}, \texttt{max\_num\_hands=1}, and \texttt{model\_complexity=1}. The neural network outputs 21 landmarks in normalized camera space $\mathbf{l}_i = (x_i, y_i, z_i) \in [0, 1]^3$, where $x$ and $y$ represent image plane coordinates and $z$ represents relative depth relative to the wrist joint.
    \item \textbf{Kinematic Feature Extraction}: The raw landmarks are ingested by the \texttt{GestureRecognizer} module, which evaluates anatomical scale baselines, joint angles, finger extension ratios, and the normalized pinch metric (formally derived in Section~\ref{sec:methodology}).
    \item \textbf{JSON Telemetry Serialization}: The resulting state is packaged into a compact JSON telemetry frame containing the discrete gesture state (\texttt{IDLE}, \texttt{OPEN\_PALM}, \texttt{FIST}, \texttt{PINCH}, \texttt{PEACE}), the toggle lock state, 2D filtered cursor coordinates, palm roll/swipe angles, the normalized pinch ratio, and high-resolution Unix timestamps. Telemetry frames are broadcast over a local WebSocket server at \texttt{ws://localhost:8000/ws/tracking}.
\end{enumerate}

\subsection{Browser CAD Client and Dual-Path Rendering}
The client-side modeling interface is implemented in React 18, utilizing Three.js and React Three Fiber (R3F) for WebGL scene graph management.

A critical design challenge in browser-based spatial interaction is handling streaming coordinate data without inducing frame drops. When high-frequency spatial streams ($30$--$60\text{ Hz}$) are routed through declarative React component state, the frequent state updates trigger continuous Virtual DOM reconciliation, leading to heap thrashing and garbage collection (GC) stutter.

To eliminate this bottleneck, GestureCAD implements a \textbf{dual-path execution architecture}:

\begin{enumerate}
    \item \textbf{High-Frequency Render Path (Bypassing Virtual DOM)}: Continuous tracking data (cursor coordinates, real-time stroke preview lines, and plane rotation handles) are maintained in mutable JavaScript references (\texttt{useRef}) rather than reactive component state. In the Three.js render loop (\texttt{useFrame}), the engine mutates the underlying WebGL \texttt{BufferGeometry.attributes.position.array} (\texttt{Float32Array}) directly in GPU-allocated memory and sets \texttt{needsUpdate = true}. This direct buffer mutation ensures continuous 60 FPS viewport feedback with zero Virtual DOM reconciliation overhead during active sketching.
    \item \textbf{Transactional CAD State Path (Command Pattern)}: Discrete geometric operations---such as committing a finished 2D profile, extruding a solid, executing a face selection, or performing Undo/Redo---are treated as discrete transactions. These operations are dispatched to a centralized Zustand state store and managed via a formal Command Pattern (\texttt{CommandManager.js}). Each command maintains inverse operations for deterministic Undo/Redo reversibility, updating the scene graph only upon transaction completion.
\end{enumerate}

\subsection{Coordinate Spaces and Transformation Pipeline}
GestureCAD coordinates geometric interactions across four distinct coordinate reference frames:
\begin{itemize}
    \item \textbf{Camera Normalized Space $\mathcal{S}_{\text{cam}}$}: Hand landmarks $\mathbf{l}_i \in [0, 1]^2$ normalized to webcam image dimensions.
    \item \textbf{Normalized Device Coordinates (NDC) $\mathcal{S}_{\text{ndc}}$}: Viewport screen coordinates mapped to $[-1, 1]^3$ for WebGL camera projection.
    \item \textbf{World Space $\mathcal{S}_{\text{world}}$}: Global three-dimensional Cartesian coordinate frame $\mathbb{R}^3$ wherein Three.js scene cameras, lighting, and instantiated CAD objects reside.
    \item \textbf{Local Working-Plane Space $\mathcal{S}_{\text{plane}}$}: 2D parametric coordinate space $(u, v) \in \mathbb{R}^2$ established on an active working plane (XY, YZ, XZ) or dynamically attached to an existing polyhedron face via an orthonormal transformation $\mathbf{R} \in SO(3)$.
\end{itemize}
The mathematical mappings connecting these coordinate spaces are detailed in the subsequent section.

\section{Methodology and Kinematic Conditioning}
\label{sec:methodology}

Transforming unconstrained monocular hand tracking into precision CAD input requires resolving physiological tremor, perspective scaling, and discrete gesture ambiguity. This section formalizes the mathematical and algorithmic pipeline governing signal conditioning, gesture state estimation, and spatial unprojection.

\subsection{Articulated Skeletal Hand Representation}
Let $\mathcal{L}_t = \{\mathbf{l}_i(t) = (x_i(t), y_i(t), z_i(t)) \in \mathbb{R}^3\}_{i=0}^{20}$ denote the set of 21 articulated skeletal hand landmarks estimated by MediaPipe at discrete time frame $t$. Key landmark indices include:
\begin{itemize}
    \item $\mathbf{l}_0$: Carpal wrist center.
    \item $\mathbf{l}_4$: Distal thumb tip.
    \item $\mathbf{l}_8$: Distal index fingertip.
    \item $\mathbf{l}_9$: Middle finger metacarpophalangeal (MCP) knuckle joint.
    \item $\mathbf{l}_{12}, \mathbf{l}_{16}, \mathbf{l}_{20}$: Tips of middle, ring, and pinky fingers.
\end{itemize}

\begin{figure}[t]
    \centering
    \fbox{%
        \parbox[c][5.4cm][c]{0.92\linewidth}{%
            \centering\small\sffamily
            \textbf{Figure 3: Gesture Interaction Workflow Placeholder}\\
            \vspace{0.3cm}
            \textit{Intended Diagram:}\\
            State transition diagram of the finite state automaton (FSM):\\
            -- States: \texttt{IDLE}, \texttt{OPEN\_PALM}, \texttt{FIST}, \texttt{PINCH}, \texttt{PEACE}.\\
            -- Transitions governed by asymmetric thresholds ($\tau_{\text{start}}=0.44$, $\tau_{\text{release}}=0.54$) and temporal debounce windows ($N_{\text{start}}=1$, $N_{\text{release}}=2$ frames).\\
            -- Lock toggle state triggered by sustained \texttt{FIST} gesture ($N \ge 5$ frames).\\
            -- Mapping of states to CAD operations: \texttt{PINCH} $\to$ Draw/Extrude; \texttt{OPEN\_PALM} $\to$ Hover/Navigate; \texttt{FIST} $\to$ Lock/Commit; \texttt{PEACE} $\to$ Plane Rotate.\\
            \vspace{0.3cm}
            \footnotesize Expected file: \texttt{paper/figures/fig3_fsm_workflow.png}
        }%
    }
    \caption{Finite state automaton governing gesture recognition and command mapping in GestureCAD. Dual-threshold hysteresis and persistence debouncing prevent limit-cycle chatter at gesture decision boundaries.}
    \label{fig:gesture_workflow}
\end{figure}

\subsection{Anatomical Scale Invariance}
In monocular vision, a user's distance from the camera $Z_{\text{cam}}$ varies naturally during interaction. Under standard pinhole camera projection with focal length $f$, a physical spatial distance $\Delta X$ at depth $Z$ projects to pixel length $\Delta x_{\text{px}}$ according to:
\begin{equation}
\Delta x_{\text{px}} = \frac{\Delta X \cdot f}{Z_{\text{cam}}}
\label{eq:pinhole}
\end{equation}
Consequently, raw Euclidean pixel distances $\|\mathbf{l}_4 - \mathbf{l}_8\|_2$ vary dramatically as the user leans forward or backward, causing fixed-pixel gesture thresholds to trigger prematurely at close distances or fail entirely at far distances.

To achieve scale invariance with respect to camera distance and individual user hand anthropometry, GestureCAD defines an anatomical normalization baseline $S_{\text{hand}}$ using the internal Euclidean distance from the wrist $\mathbf{l}_0$ to the middle MCP knuckle $\mathbf{l}_9$:
\begin{equation}
S_{\text{hand}}(t) = \|\mathbf{l}_9(t) - \mathbf{l}_0(t)\|_2 = \sqrt{\sum_{k \in \{x, y, z\}} (l_{9, k}(t) - l_{0, k}(t))^2}
\label{eq:hand_size}
\end{equation}
The middle MCP joint is biomechanically optimal as a reference anchor because the carpal-metacarpal distance remains virtually rigid during finger articulation (unlike finger tips, which translate during curling or grasping).

The dimensionless, scale-invariant pinch metric $d_{\text{pinch}}(t)$ is then defined as:
\begin{equation}
d_{\text{pinch}}(t) = \frac{\|\mathbf{l}_4(t) - \mathbf{l}_8(t)\|_2}{S_{\text{hand}}(t)}
\label{eq:d_pinch}
\end{equation}
\textit{Analytical Scale Invariance Property}: Under ideal pinhole perspective projection without out-of-plane rotation, substituting Eq.~\eqref{eq:pinhole} into Eq.~\eqref{eq:d_pinch} yields:
\begin{equation}
d_{\text{pinch}}(t) = \frac{\frac{\Delta X_{\text{pinch}} \cdot f}{Z_{\text{cam}}}}{\frac{\Delta X_{\text{baseline}} \cdot f}{Z_{\text{cam}}}} = \frac{\Delta X_{\text{pinch}}}{\Delta X_{\text{baseline}}}
\label{eq:scale_invariance_proof}
\end{equation}
Because the depth $Z_{\text{cam}}$ and focal length $f$ cancel identically, $d_{\text{pinch}}$ is analytically invariant to camera distance ($0.0\%$ theoretical coefficient of variation across depths). 

\textit{Real-World Qualification}: In physical camera interaction, extreme out-of-plane hand tilt induces differential perspective foreshortening between non-coplanar finger segments, introducing residual variance. However, within typical interaction cones ($\pm 25^\circ$ hand pitch and yaw), the metric provides robust depth independence.

\subsection{Dual-Threshold Hysteresis State Automaton}
When an interaction threshold is governed by a single boundary $\tau$, high-frequency physiological tremor near $\tau$ triggers rapid, uncontrolled limit-cycle transitions between states. GestureCAD suppresses this boundary chatter using a dual-threshold Schmitt trigger automaton with temporal persistence debouncing.

Let $\sigma_t \in \{\text{IDLE}, \text{PINCH}\}$ represent the binary activation state. Transitions are defined by asymmetric activation and release boundaries:
\begin{equation}
\sigma_t = \begin{cases}
\text{PINCH}, & \text{if } \sigma_{t-1} = \text{IDLE} \land d_{\text{pinch}}(t) < \tau_{\text{start}} \land C_{\text{debounce}} \ge N_{\text{start}} \\
\text{IDLE}, & \text{if } \sigma_{t-1} = \text{PINCH} \land d_{\text{pinch}}(t) > \tau_{\text{release}} \land C_{\text{debounce}} \ge N_{\text{release}} \\
\sigma_{t-1}, & \text{otherwise}
\end{cases}
\label{eq:hysteresis}
\end{equation}
where $\tau_{\text{start}} = 0.44$, $\tau_{\text{release}} = 0.54$, and $C_{\text{debounce}}$ is an integer counter requiring state persistence over $N_{\text{start}} = 1$ and $N_{\text{release}} = 2$ consecutive frames.

The full gesture recognizer evaluates five discrete poses:
\begin{enumerate}
    \item \textbf{PINCH}: Index and thumb tips meet ($d_{\text{pinch}} < \tau_{\text{start}}$) with ring and pinky fingers relaxed (extension ratios $> 0.80$). Used for continuous drawing, profile vertex placement, and 3D extrusion pulling.
    \item \textbf{OPEN\_PALM}: All four non-thumb fingers fully extended (extension ratios $> 1.35$). Used for passive cursor hovering, plane translation, and face inspection.
    \item \textbf{FIST}: All four fingers folded toward the palm (extension ratios $< 0.85$). When held stably for $N \ge 5$ frames, it toggles a safety lock state, preventing accidental drawing activations.
    \item \textbf{PEACE}: Index and middle fingers extended, ring and pinky folded, and thumb separated ($d_{\text{pinch}} > 0.50$). Used for angular plane reorientation and viewport rotation.
    \item \textbf{IDLE}: Default resting state when no structured pose criteria are met.
\end{enumerate}

\subsection{Multi-Stage Kinematic Tremor Attenuation}
Raw landmark position streams contain human micro-tremor and optical tracking jitter: $\tilde{\mathbf{p}}_t = \mathbf{p}_t + \boldsymbol{\epsilon}_t$. Applying standard fixed-parameter low-pass filtering introduces noticeable phase lag during rapid hand sweeps while failing to eliminate tremor during stationary aiming.

GestureCAD conditions the 2D cursor trajectory through a three-stage sequential filtering pipeline:

\textbf{Stage 1: Micro-Tremor Dead-Zone Gating.} Involuntary human hand tremor while attempting to hold a steady spatial position typically spans displacements below $3\text{ mm}$. We reject tiny inter-frame perturbations using a squared Euclidean dead-zone threshold:
\begin{equation}
\hat{\mathbf{p}}_t = \begin{cases}
\mathbf{p}_{t-1}, & \text{if } \|\tilde{\mathbf{p}}_t - \mathbf{p}_{t-1}\|_2^2 < \epsilon_{\text{dead}}^2 \\
\tilde{\mathbf{p}}_t, & \text{otherwise}
\end{cases}
\label{eq:deadzone}
\end{equation}
where $\epsilon_{\text{dead}} = 0.003\text{ m}$ ($3\text{ mm}$ in normalized hand coordinate space).

\textbf{Stage 2: Velocity-Clamped Outlier Rejection.} To mitigate momentary optical tracking glitches caused by self-occlusion or lighting shifts, maximum inter-frame displacement is bounded by a kinematic velocity cap $\delta_{\max}$:
\begin{equation}
\mathbf{p}^*_t = \mathbf{p}_{t-1} + \min\left(1.0, \, \frac{\delta_{\max}}{\|\hat{\mathbf{p}}_t - \mathbf{p}_{t-1}\|_2}\right) (\hat{\mathbf{p}}_t - \mathbf{p}_{t-1})
\label{eq:velocity_clamp}
\end{equation}
where $\delta_{\max} = 2.5\text{ m/frame}$.

\textbf{Stage 3: Velocity-Adaptive Exponential Moving Average (EMA).} The smoothing weight $\alpha(v_t)$ dynamically adapts as a monotonic function of instantaneous hand speed $v_t = \|\mathbf{p}^*_t - \mathbf{p}_{t-1}\|_2$:
\begin{equation}
\alpha(v_t) = \min\left(1.0, \, \alpha_0 + \min\left(\alpha_{\text{boost}}, \, k_v \cdot v_t\right)\right)
\label{eq:adaptive_alpha}
\end{equation}
\begin{equation}
\mathbf{p}_t = \mathbf{p}_{t-1} + \alpha(v_t) \cdot (\mathbf{p}^*_t - \mathbf{p}_{t-1})
\label{eq:ema}
\end{equation}
with baseline smoothing $\alpha_0 = 0.35$, dynamic boost limit $\alpha_{\text{boost}} = 0.25$, and velocity scaling coefficient $k_v = 0.50$. When the hand is moving slowly for fine positioning, $\alpha \to \alpha_0$, providing heavy low-pass filtering. During rapid strokes, $\alpha \to 0.60$--$1.0$, minimizing trailing latency.

\begin{figure}[t]
    \centering
    \fbox{%
        \parbox[c][5.4cm][c]{0.92\linewidth}{%
            \centering\small\sffamily
            \textbf{Figure 4: Working Planes and Unprojection Placeholder}\\
            \vspace{0.3cm}
            \textit{Intended Screenshot:}\\
            Visual demonstration of parametric working plane interaction in 3D viewport:\\
            -- Viewport showing primary coordinate planes: $XY$ (Ground), $YZ$ (Side), and $XZ$ (Front);\\
            -- The visual ray $\vec{r}(t) = \vec{r}_0 + t \vec{d}$ projecting from the virtual camera through the NDC cursor onto the active grid plane;\\
            -- The calculated 3D intersection hit-point $\vec{p}_{\text{hit}}$ and its mapping into 2D local in-plane coordinates $(u, v)$;\\
            -- Visual snapping indicators highlighting vertex, midpoint, and edge snap targets.\\
            \vspace{0.3cm}
            \footnotesize Expected file: \texttt{paper/figures/fig4_working_planes.png}
        }%
    }
    \caption{Parametric working planes and perspective unprojection. Monocular hand cursor coordinates are mapped to a 3D camera ray and intersected with the active drawing plane, establishing an orthonormal $(u, v)$ coordinate system for 2D profile authoring.}
    \label{fig:working_planes}
\end{figure}

\subsection{Perspective Unprojection and Working-Plane Raycasting}
To author 3D geometry from 2D gesture coordinates, the system projects screen cursor positions onto active 3D parametric working planes (Figure~\ref{fig:working_planes}).

A normalized screen coordinate $(x, y) \in [0, 1]^2$ is transformed into Normalized Device Coordinates (NDC) $\mathbf{x}_{\text{ndc}} \in [-1, 1]^3$:
\begin{equation}
\mathbf{x}_{\text{ndc}} = \begin{bmatrix} 2x - 1 \\ -(2y - 1) \\ 1 \end{bmatrix}
\label{eq:ndc}
\end{equation}
Using the camera world position $\mathbf{C}_{\text{camera}} \in \mathbb{R}^3$ and the inverse of the view-projection matrix $\mathbf{M}_{\text{vp}}^{-1} = (\mathbf{P}_{\text{proj}} \mathbf{V}_{\text{view}})^{-1}$, the world-space ray origin $\vec{r}_0$ and unit direction vector $\vec{d}$ are evaluated:
\begin{equation}
\vec{r}_0 = \mathbf{C}_{\text{camera}}, \quad \vec{d} = \frac{\mathbf{M}_{\text{vp}}^{-1} \mathbf{x}_{\text{ndc}}}{\|\mathbf{M}_{\text{vp}}^{-1} \mathbf{x}_{\text{ndc}}\|_2}
\label{eq:ray}
\end{equation}

An active planar working plane $\Pi$ is parameterized by a reference center point $\vec{p}_0 \in \mathbb{R}^3$ and unit outward normal $\vec{n} \in \mathbb{R}^3$. The ray-plane intersection parameter $t^*$ is given by:
\begin{equation}
t^* = \frac{(\vec{p}_0 - \vec{r}_0) \cdot \vec{n}}{\vec{d} \cdot \vec{n}}, \quad \text{valid for } |\vec{d} \cdot \vec{n}| > 10^{-6} \land t^* \ge 0
\label{eq:ray_plane_intersection}
\end{equation}
The resulting 3D hit point $\vec{p}_{\text{hit}} = \vec{r}_0 + t^* \vec{d}$ is projected onto local 2D working plane coordinates $(u, v)$ via in-plane tangent and bitangent unit vectors $(\vec{u}_{\text{tangent}}, \vec{v}_{\text{bitangent}})$:
\begin{equation}
u = (\vec{p}_{\text{hit}} - \vec{p}_0) \cdot \vec{u}_{\text{tangent}}, \quad v = (\vec{p}_{\text{hit}} - \vec{p}_0) \cdot \vec{v}_{\text{bitangent}}
\label{eq:plane_uv}
\end{equation}

\subsection{Geometric Snapping and Shape Regularization}
Because mid-air hand sketching lacks physical friction, users cannot easily close polygons or align vertices manually. GestureCAD incorporates two assisting engines:

\textbf{1. 3D Geometric Snapping}: When placing a point $\mathbf{p}$, the snapping engine searches for existing geometric candidates within a screen-space tolerance radius $r_{\text{snap}} = 0.25\text{ m}$:
\begin{itemize}
    \item \textbf{Vertex Snapping}: Pulls $\mathbf{p}$ to existing vertex coordinates $\|\mathbf{p} - \mathbf{v}_k\| < r_{\text{snap}}$, guaranteeing exact numerical profile closure.
    \item \textbf{Midpoint Snapping}: Snaps $\mathbf{p}$ to edge midpoints $\frac{1}{2}(\mathbf{v}_i + \mathbf{v}_j)$.
    \item \textbf{Grid Snapping}: Quantizes $\mathbf{p}$ to increments of $0.5\text{ m}$ on the active working plane.
\end{itemize}

\textbf{2. Heuristic Shape Regularization}: A continuous freehand stroke $\mathcal{S} = \{\mathbf{q}_k\}_{k=0}^{M-1}$ is evaluated upon pen release. The regularizer evaluates stroke perimeter, aspect ratio, end-to-end closure distance $\|\mathbf{q}_{M-1} - \mathbf{q}_0\|$, and deviation from ideal geometric primitives:
\begin{itemize}
    \item If $\|\mathbf{q}_{M-1} - \mathbf{q}_0\| / L_{\text{perimeter}} < 0.15$ and the bounding box aspect ratio $\approx 1.0$, the stroke is regularized into a circle or rectangle.
    \item If the stroke has 3 prominent directional inflection points, it is regularized into an equilateral or right-angled triangle.
    \item Open strokes with collinearity index $> 0.95$ are simplified to straight 2-point line segments.
\end{itemize}
This regularization ensures clean, closed 2D profiles suitable for downstream solid generation.

\section{CAD Geometry Generation and Validation}
\label{sec:geometry}

A core distinction between artistic spatial sketching and engineering CAD is the mathematical rigor of the underlying geometric model. GestureCAD implements a custom polyhedral solid modeling kernel that generates closed, watertight surface meshes via linear extrusion and rotational sweeping, coupled with automated topological validation.

\subsection{2D Planar Profile Parameterization and Winding Orientation}
A 2D sketch profile is parameterized as an ordered sequence of planar vertices in local working-plane coordinates: $\mathcal{P} = \{\mathbf{q}_i = (u_i, 0, v_i)\}_{i=0}^{N-1}$. 

To construct consistent outward-facing normals during subsequent 3D sweeping, the profile must possess a deterministic counter-clockwise (CCW) winding order. The signed planar area $A(\mathcal{P})$ is evaluated using the Gauss-Green / Shoelace formula:
\begin{equation}
A(\mathcal{P}) = \frac{1}{2} \sum_{i=0}^{N-1} (u_i v_{i+1} - u_{i+1} v_i), \quad \text{with } \mathbf{q}_N \equiv \mathbf{q}_0
\label{eq:shoelace}
\end{equation}
If $A(\mathcal{P}) < 0$, the profile vertices traverse clockwise (CW); the order is reversed: $\mathcal{P} \leftarrow \text{reverse}(\mathcal{P})$, guaranteeing $A(\mathcal{P}) > 0$.

\begin{figure}[t]
    \centering
    \fbox{%
        \parbox[c][5.4cm][c]{0.92\linewidth}{%
            \centering\small\sffamily
            \textbf{Figure 5: Sketch-to-Solid Pipeline Placeholder}\\
            \vspace{0.3cm}
            \textit{Intended Sequential Progression (Horizontal Strip):}\\
            (a) \textbf{Planar Profile}: 2D closed polygon sketched on the active working plane with vertex closure snapping;\\
            (b) \textbf{Normal Sweeping}: Continuous upward displacement along plane normal $\vec{n}_{\text{plane}}$ driven by Z-Hold hand gesture;\\
            (c) \textbf{Watertight Extrusion}: Resulting closed 3D polyhedral prism with outward side quads and planar end caps;\\
            (d) \textbf{Revolved Solid}: 2D cross-section revolved $360^\circ$ around an in-plane axis via Rodrigues' formula with welded axis vertices.\\
            \vspace{0.3cm}
            \footnotesize Expected file: \texttt{paper/figures/fig5_sketch_to_solid.png}
        }%
    }
    \caption{The sketch-to-solid modeling pipeline in GestureCAD. Closed 2D profiles are swept into closed 3D polyhedra through linear extrusion or rotational sweeping.}
    \label{fig:sketch_to_solid}
\end{figure}

\subsection{Linear Extrusion and Cap Triangulation}
Linear extrusion sweeps the planar profile $\mathcal{P}$ along the working plane normal vector $\vec{n}_{\text{plane}}$ by an extrusion distance $h > 0$:
\begin{equation}
\Omega = \mathcal{P} \times [0, h] \cdot \vec{n}_{\text{plane}}
\label{eq:extrusion_sweep}
\end{equation}
The extrusion kernel constructs:
\begin{enumerate}
    \item \textbf{Bottom Ring Vertices}: $\mathbf{v}_i = [u_i, 0, v_i]$ for $i \in \{0, \dots, N-1\}$.
    \item \textbf{Top Ring Vertices}: $\mathbf{v}_{i+N} = [u_i, h, v_i]$ for $i \in \{0, \dots, N-1\}$. Total vertex count: $V = 2N$.
    \item \textbf{Side Wall Faces}: For each segment between vertex $i$ and $j = (i+1) \bmod N$, two counter-clockwise outward-oriented triangles are generated:
    \begin{equation}
    \mathcal{F}_{\text{side}, 1} = [i, i+N, j+N], \quad \mathcal{F}_{\text{side}, 2} = [i, j+N, j]
    \label{eq:side_quads}
    \end{equation}
    yielding $2N$ triangular side faces.
    \item \textbf{Cap Triangulation}: In the current implementation, planar end caps are constructed using a triangle fan rooted at index 0 for the bottom cap and index $N$ for the top cap:
    \begin{align}
    \mathcal{F}_{\text{bottom}, i} &= [0, i, i+1], & i \in \{1, \dots, N-2\} \label{eq:bottom_cap} \\
    \mathcal{F}_{\text{top}, i} &= [N, N+i+1, N+i], & i \in \{1, \dots, N-2\} \label{eq:top_cap}
    \end{align}
    yielding $N-2$ triangles per cap, or $2N - 4$ cap faces in total.
\end{enumerate}

\subsubsection{Critical Computational Geometry Audit of Fan Triangulation}
A rigorous computational geometry analysis reveals an important mathematical distinction between topological graph connectivity and spatial geometric embedding:
\begin{itemize}
    \item \textbf{Convex Profiles}: For any strictly convex polygon (e.g., rectangles, circles, regular polygons), the triangle fan from vertex 0 produces a valid, non-overlapping planar triangulation that fully tiles the interior.
    \item \textbf{Concave Profiles}: For non-convex profiles (e.g., L-shapes, T-shapes, star polygons), chords connecting vertex 0 to vertex $i$ may lie \textit{outside} the polygon boundary or cross other edges. Consequently, fan cap triangles overlap and self-intersect in $\mathbb{R}^3$.
    \item \textbf{Topological vs. Geometric Validity}: Interestingly, because the half-edge graph topology remains isomorphic to a triangulated sphere, the mesh still reports an Euler characteristic of $\chi = 2$ and 0 boundary edges. Furthermore, under Gauss's Divergence Theorem, overlapping exterior triangles traverse in reverse orientation, causing their signed volume contributions to cancel out algebraically and still match the analytical volume. However, \textit{physically}, the resulting surface mesh is an immersed 2-manifold with self-intersections rather than an embedded manifold, causing rendering artifacts and slicing errors in CAM software.
    \item \textbf{Required Generalization}: To support arbitrary simple concave polygons without self-intersection, cap generation requires Ear-Clipping (e.g., Earcut~\cite{held2001vrp}) or constrained Delaunay triangulation. (Notably, GestureCAD's \texttt{RevolveEngine} already uses Earcut for partial end caps, as detailed below).
\end{itemize}

\subsection{Rotational Sweeping (Revolve Engine)}
The \texttt{RevolveEngine} sweeps a 2D cross-sectional profile $\mathcal{P}$ around an arbitrary in-plane axis of revolution defined by point $\mathbf{a}_0$ and unit direction $\vec{v}$.

For each profile vertex $\mathbf{p} \in \mathcal{P}$, the relative vector $\mathbf{r} = \mathbf{p} - \mathbf{a}_0$ is decomposed into components parallel and perpendicular to $\vec{v}$:
\begin{equation}
\mathbf{r}_{\parallel} = (\mathbf{r} \cdot \vec{v}) \vec{v}, \quad \mathbf{r}_{\perp} = \mathbf{r} - \mathbf{r}_{\parallel}
\label{eq:revolve_decomp}
\end{equation}
The rotation of $\mathbf{p}$ through angle $\phi$ is computed via Rodrigues' rotation formula:
\begin{equation}
\mathbf{p}_{\text{rot}}(\phi) = \mathbf{a}_0 + \mathbf{r}_{\parallel} + \cos(\phi) \mathbf{r}_{\perp} + \sin(\phi) (\vec{v} \times \mathbf{r}_{\perp})
\label{eq:rodrigues}
\end{equation}

The rotational sweep incorporates two critical solid modeling safeguards:
\begin{enumerate}
    \item \textbf{On-Axis Vertex Welding}: If a profile vertex lies on the rotation axis ($\|\mathbf{r}_{\perp}\| < 10^{-4}\text{ m}$), revolving it would create degenerate zero-area quads. The engine identifies on-axis vertices and assigns them a single fixed vertex index across all angular rings. Adjacent faces collapse into fan triangles terminating at the axis.
    \item \textbf{Seam Stitching and Partial Caps}: For a complete $360^\circ$ revolution, ring $M$ wraps seamlessly to ring 0 without duplicate vertices. For partial revolutions ($\theta < 360^\circ$), the start and end boundary profiles are sealed using Earcut triangulation (\texttt{THREE.ShapeUtils.triangulateShape}), guaranteeing valid non-intersecting end caps for arbitrary profiles.
\end{enumerate}

\begin{figure}[t]
    \centering
    \fbox{%
        \parbox[c][5.4cm][c]{0.92\linewidth}{%
            \centering\small\sffamily
            \textbf{Figure 6: Geometry and Topology Validation Placeholder}\\
            \vspace{0.3cm}
            \textit{Intended Screenshot:}\\
            Inspection of a generated polyhedral solid mesh showing:\\
            -- Viewport wireframe showing 2-manifold edge connectivity with no unshared boundaries;\\
            -- Visual arrows displaying computed outward-pointing unit face normal vectors $\vec{n}_k$;\\
            -- The interactive CAD Inspector overlay verifying topological invariants:\\
            * Euler characteristic: $\chi = V - E + F = 2.000$;\\
            * Boundary edges: $\partial \Omega = 0$;\\
            * Non-manifold edges: $0$;\\
            * Exact Divergence Theorem volume: $V_{\text{div}} = 24.375\text{ m}^3$ matching analytical volume.\\
            \vspace{0.3cm}
            \footnotesize Expected file: \texttt{paper/figures/fig6_mesh_validation.png}
        }%
    }
    \caption{Topological invariant validation in GestureCAD. Every generated solid is verified for watertightness ($\partial \Omega = 0$), 2-manifold edge degree ($d(e) = 2$), orientation consistency, Euler characteristic ($\chi = 2$), and Divergence Theorem volume.}
    \label{fig:mesh_validation}
\end{figure}

\subsection{Topological Invariant Validation Suite}
Every geometric solid instantiated in GestureCAD is represented by a \texttt{SolidMesh} data structure and subjected to an automated topological validation suite (Figure~\ref{fig:mesh_validation}):

\textbf{1. Edge Adjacency and 2-Manifold Verification}: The engine constructs a half-edge adjacency map. Each undirected edge $e = (v_a, v_b)$ records the list of incident faces and their traversal directions. A valid 2-manifold solid must satisfy:
\begin{equation}
|\mathcal{F}(e)| = 2, \quad \forall e \in \mathcal{E}
\label{eq:manifold_edge}
\end{equation}
Edges with $|\mathcal{F}(e)| = 1$ represent open boundary tears ($\partial \Omega \neq \emptyset$), indicating a non-watertight mesh. Edges with $|\mathcal{F}(e)| > 2$ represent non-manifold junctions where three or more surface sheets intersect.

\textbf{2. Normal Orientation Consistency}: For every edge shared by two faces $\mathcal{F}_1$ and $\mathcal{F}_2$, the directed edge traversals must be strictly opposed:
\begin{equation}
\text{direction}(\mathcal{F}_1, e) + \text{direction}(\mathcal{F}_2, e) = 0
\label{eq:orientation_consistency}
\end{equation}
ensuring all surface normals point consistently outward.

\textbf{3. Euler Characteristic Invariant}: For any closed, orientable, simply connected 2-manifold polyhedron of genus $g = 0$ (homeomorphic to a 2-sphere), the Euler characteristic $\chi$ is an invariant topological integer:
\begin{equation}
\chi = V - E + F = 2(1 - g) = 2
\label{eq:euler}
\end{equation}
For an extruded $N$-gon mesh, $V = 2N$, $F = 4N - 4$, and $E = 6N - 6$. Substituting these counts into Eq.~\eqref{eq:euler}:
\begin{equation}
\chi = (2N) - (6N - 6) + (4N - 4) = 2
\label{eq:euler_proof}
\end{equation}
\textit{Topological Scope}: The invariant $\chi = 2$ strictly characterizes genus-0 solids. Prisms containing through-holes or hollow toroidal cavities possess genus $g \ge 1$ ($\chi \le 0$) and are outside the current kernel's topological domain.

\textbf{4. Exact Volumetric Integration via Divergence Theorem}: The signed volume $V(\Omega)$ of the enclosed solid domain $\Omega$ is computed directly from its boundary surface $\partial \Omega$ without voxelization. Applying Gauss's Divergence Theorem to the vector field $\vec{F}(\mathbf{x}) = \frac{1}{3}\mathbf{x}$ (where $\nabla \cdot \vec{F} = 1$):
\begin{equation}
V(\Omega) = \iiint_{\Omega} (\nabla \cdot \vec{F}) \, dV = \iint_{\partial \Omega} (\vec{F} \cdot \vec{n}) \, dA
\label{eq:divergence_theorem}
\end{equation}
For a closed surface composed of planar triangular facets $\mathcal{F}_k = (\mathbf{a}_k, \mathbf{b}_k, \mathbf{c}_k)$ with outward CCW winding, this reduces to the discrete signed tetrahedral sum:
\begin{equation}
V(\Omega) = \frac{1}{6} \sum_{k=1}^{F} \mathbf{a}_k \cdot (\mathbf{b}_k \times \mathbf{c}_k)
\label{eq:divergence_volume}
\end{equation}

\begin{figure}[t]
    \centering
    \fbox{%
        \parbox[c][5.4cm][c]{0.92\linewidth}{%
            \centering\small\sffamily
            \textbf{Figure 7: Multi-Surface Sketching Workflow Placeholder}\\
            \vspace{0.3cm}
            \textit{Intended Screenshot:}\\
            Demonstration of dynamic sketch-on-face plane attachment:\\
            -- A base extruded rectangular prism with named faces (Top, Front, Right);\\
            -- The user raycasts and selects the top face $\to$ FaceEnumerator constructs local orthonormal frame $\mathbf{R}_{\text{face}} \in SO(3)$;\\
            -- A local 2D sketch plane is attached directly onto the selected face;\\
            -- A secondary circular profile is sketched on the face and extruded upward into an aligned cylindrical boss feature.\\
            \vspace{0.3cm}
            \footnotesize Expected file: \texttt{paper/figures/fig7_sketch_on_face.png}
        }%
    }
    \caption{Multi-surface sketching in GestureCAD. The \texttt{FaceEnumerator} derives a local orthonormal coordinate basis $\mathbf{R} \in SO(3)$ from selected polyhedron faces, permitting hierarchical sketch-on-face operations directly in 3D space.}
    \label{fig:sketch_on_face}
\end{figure}

\subsection{Multi-Surface Decomposition and Local $SO(3)$ Tangent Frames}
A defining capability of professional CAD is the capacity to sketch new features directly on existing solid faces (Figure~\ref{fig:sketch_on_face}). GestureCAD implements this via the \texttt{FaceEnumerator} pipeline:

When a user raycasts onto a solid mesh with outward face normal $\vec{n} \in \mathbb{R}^3$, the system constructs a local, right-handed orthonormal coordinate basis $\mathbf{R}_{\text{face}} \in SO(3)$ on the surface:
\begin{enumerate}
    \item To avoid numerical gimbal lock when $\vec{n}$ aligns with the global vertical axis, a reference vector $\vec{r}_{\text{ref}}$ is selected:
    \begin{equation}
    \vec{r}_{\text{ref}} = \begin{cases}
    [0, 1, 0]^T, & \text{if } |\vec{n} \cdot [0, 1, 0]^T| < 0.9 \\
    [1, 0, 0]^T, & \text{otherwise}
    \end{cases}
    \label{eq:ref_vector}
    \end{equation}
    \item In-plane unit tangent $\vec{t}$ and bitangent $\vec{b}$ vectors are computed via Gram-Schmidt cross products:
    \begin{equation}
    \vec{t} = \frac{\vec{n} \times \vec{r}_{\text{ref}}}{\|\vec{n} \times \vec{r}_{\text{ref}}\|_2}, \quad \vec{b} = \vec{t} \times \vec{n}
    \label{eq:tangent_basis}
    \end{equation}
    \item The face-attached rotation matrix $\mathbf{R}_{\text{face}} = [\vec{t} \;\; \vec{n} \;\; \vec{b}]$ satisfies:
    \begin{equation}
    \det(\mathbf{R}_{\text{face}}) = +1, \quad \mathbf{R}_{\text{face}}^T \mathbf{R}_{\text{face}} = \mathbf{I}
    \label{eq:so3_property}
    \end{equation}
\end{enumerate}
For curved cylindrical surfaces, the normal $\vec{n}$ is computed radially from the longitudinal axis to the raycast hit point, establishing a local tangent plane for surface-attached sketching.

\section{Experimental Evaluation and Results}
\label{sec:experiments}

To ensure complete academic transparency, all experimental evaluations in GestureCAD are explicitly categorized by scientific provenance:
\begin{itemize}
    \item \textbf{REAL}: Evaluated on real executing software code or physical hardware.
    \item \textbf{SIMULATED (Parameter Sweeps)}: Evaluated across multi-seed stochastic parameter sweeps with reported 95\% confidence intervals (CIs) and min-max ranges.
    \item \textbf{ANALYTIC}: Derived from exact mathematical and geometric theorems.
    \item \textbf{UNMEASURED}: Explicitly documented as unmeasured pending specialized physical laboratory apparatus or human trials.
\end{itemize}
Table~\ref{tab:provenance} summarizes the evaluation matrix.

\begin{table}[t]
\caption{Scientific provenance and methodology matrix.}
\label{tab:provenance}
\centering
\small
\begin{tabular}{llcc}
\toprule
\textbf{ID} & \textbf{Evaluated Dimension} & \textbf{Provenance} & \textbf{Sample Scope} \\
\midrule
B1 & Software CAD Subsystem Latency & \textbf{REAL} & 5,000 iterations \\
B2 & Physical Motion-to-Photon Latency & \textbf{UNMEASURED} & Requires high-speed rig \\
B3 & Cursor Jitter Attenuation & \textbf{SIMULATED} & 180 runs (90k pts) \\
B4 & Hysteresis State Chatter Reduction & \textbf{SIMULATED} & 180 runs (450k frames) \\
B6 & Distance Scale Invariance & \textbf{ANALYTIC} & Pinhole projection model \\
B9 & Topology \& Divergence Volume & \textbf{REAL + ANALYTIC} & 20 polyhedra \\
B10 & Shape Regularization Accuracy & \textbf{SIMULATED} & 180 runs (18k strokes) \\
B11 & Snapping Targeting Precision & \textbf{SIMULATED} & 180 runs (36k queries) \\
RL-1 & Webcam Capture Throughput & \textbf{REAL} & 141 frames physical cam \\
\bottomrule
\end{tabular}
\end{table}

\begin{figure*}[t]
    \centering
    \fbox{%
        \parbox[c][5.4cm][c]{0.92\linewidth}{%
            \centering\small\sffamily
            \textbf{Figure 8: Experimental Benchmark Results Placeholder}\\
            \vspace{0.3cm}
            \textit{Intended Dual-Panel Benchmark Plots:}\\
            (a) \textbf{Left Panel -- Tremor Attenuation Sensitivity Sweep (Benchmark B3)}: Plot showing RMS cursor jitter (mm) as a function of input noise level ($\sigma = 1.0$ to $6.0\text{ mm}$), comparing raw unconditioned input against GestureCAD filtered output across 180 simulation runs with 95\% confidence interval bands;\\
            (b) \textbf{Right Panel -- Hysteresis State Stability Sweep (Benchmark B4)}: Plot comparing number of spurious gesture state flips per 2,500 boundary frames between a naive single-threshold classifier ($\tau=0.48$) and the GestureCAD dual-threshold hysteresis automaton across six decision-boundary noise levels ($\sigma = 0.015$ to $0.090$).\\
            \vspace{0.3cm}
            \footnotesize Expected file: \texttt{paper/figures/fig8_experimental_results.png}
        }%
    }
    \caption{Empirical benchmark results. (a) RMS cursor jitter reduction across increasing noise levels ($\sigma = 1.0$--$6.0\text{ mm}$); (b) Spurious state-flip suppression of the dual-threshold hysteresis automaton near the pinch decision boundary.}
    \label{fig:experimental_results}
\end{figure*}

\subsection{Software CAD Pipeline Latency (Benchmark B1)}
We evaluated the computational execution time of each discrete software CAD stage in JavaScript/Node.js over $5\,000$ consecutive iterations on an Intel Core i7-12700H CPU ($2.30\text{ GHz}$). Table~\ref{tab:latency} provides the detailed statistical breakdown.

\begin{table}[t]
\caption{Software CAD pipeline subsystem execution latency ($N = 5\,000$ iterations, values in milliseconds).}
\label{tab:latency}
\centering
\small
\begin{tabular}{lcccc}
\toprule
\textbf{Subsystem Stage} & \textbf{Mean} & \textbf{Med.} & \textbf{p99} & \textbf{Std} \\
\midrule
Ray-Plane Unprojection & 0.0045 & 0.0032 & 0.0219 & 0.0171 \\
Snapping Inference & 0.0024 & 0.0015 & 0.0068 & 0.0107 \\
Shape Regularization & 0.0836 & 0.0598 & 0.3769 & 0.1035 \\
Extrusion \& Topology Validation & 0.1009 & 0.0826 & 0.4573 & 0.0803 \\
Face Enumeration \& $SO(3)$ Basis & 0.0080 & 0.0041 & 0.0324 & 0.0279 \\
\midrule
\textbf{Cumulative Software CAD} & \textbf{0.1994} & \textbf{0.1512} & \textbf{0.8953} & \textbf{0.2395} \\
\bottomrule
\end{tabular}
\end{table}

\textbf{Finding}: The cumulative execution time of all CAD computational stages averages $0.1994\text{ ms}$ per frame ($0.1512\text{ ms}$ median, $0.8953\text{ ms}$ 99th percentile), consuming only $1.2\%$ of the standard $16.6\text{ ms}$ ($60\text{ FPS}$) rendering budget.

\subsubsection{Critical Distinction: Software Latency vs. Motion-to-Photon Latency}
We emphasize that the $0.1994\text{ ms}$ metric reflects \textit{isolated software CPU execution time} within the CAD pipeline. It does \textbf{not} represent end-to-end physical motion-to-photon latency. The full physical pipeline encompasses:
\begin{enumerate}
    \item Webcam optical exposure and frame integration ($\approx 33.3\text{ ms}$ at 30 FPS).
    \item USB hardware transfer and OS frame decode ($\approx 5$--$10\text{ ms}$).
    \item MediaPipe neural network landmark regression ($\approx 15$--$25\text{ ms}$).
    \item Inter-process WebSocket JSON serialization and transport ($\approx 1$--$3\text{ ms}$).
    \item Software CAD evaluation ($\approx 0.2\text{ ms}$, measured).
    \item Three.js WebGL GPU command compilation and rasterization ($\approx 2$--$5\text{ ms}$).
    \item Monitor display vertical sync and scanout ($\approx 8.3$--$16.6\text{ ms}$).
\end{enumerate}
Total physical motion-to-photon latency is estimated to lie between $65$ and $95\text{ ms}$. Accurately measuring this requires an external optical setup (e.g., a 240 FPS camera capturing a physical LED trigger on the finger simultaneously with display pixels). Until such specialized instrumentation is deployed, physical end-to-end latency remains formally unmeasured.

\subsection{Tremor Attenuation Sensitivity Sweep (Benchmark B3)}
We evaluated the multi-stage tremor filter across 6 input noise levels ($\sigma \in [1.0, 2.0, 3.0, 4.0, 5.0, 6.0]\text{ mm}$) and 30 random seeds ($180$ simulation runs, $90\,000$ points total in \texttt{summary\_b3\_sweep.json}). Table~\ref{tab:jitter_sweep} reports RMS jitter before and after filtering.

\begin{table}[t]
\caption{Tremor filter jitter attenuation sweep across 6 noise levels ($30$ seeds each, $90\,000$ samples total).}
\label{tab:jitter_sweep}
\centering
\small
\begin{tabular}{ccccc}
\toprule
\textbf{Noise $\sigma$} & \textbf{Raw RMS (mm)} & \textbf{Filtered (mm)} & \textbf{Reduction} & \textbf{95\% CI} \\
\midrule
$1.0\text{ mm}$ & $1.060 \pm 0.030$ & $0.573 \pm 0.013$ & $45.92\%$ & $[45.74\%, 46.10\%]$ \\
$2.0\text{ mm}$ & $2.119 \pm 0.061$ & $1.026 \pm 0.023$ & $51.58\%$ & $[51.42\%, 51.74\%]$ \\
$3.0\text{ mm}$ & $3.179 \pm 0.091$ & $1.493 \pm 0.034$ & $53.04\%$ & $[52.88\%, 53.20\%]$ \\
$4.0\text{ mm}$ & $4.238 \pm 0.122$ & $1.967 \pm 0.046$ & $53.58\%$ & $[53.42\%, 53.74\%]$ \\
$5.0\text{ mm}$ & $5.298 \pm 0.152$ & $2.446 \pm 0.057$ & $53.83\%$ & $[53.67\%, 53.99\%]$ \\
$6.0\text{ mm}$ & $6.358 \pm 0.183$ & $2.927 \pm 0.069$ & $53.96\%$ & $[53.79\%, 54.13\%]$ \\
\midrule
\textbf{Grand Mean} & — & — & \textbf{52.43\%} & \textbf{[52.16\%, 52.71\%]} \\
\bottomrule
\end{tabular}
\end{table}

\textbf{Finding}: Across all 180 runs, the filter achieves a grand mean RMS jitter reduction of $52.43\%$ [$95\%$ CI: $52.16\%$–$52.71\%$, range: $45.0\%$–$58.6\%$]. The attenuation scales from $45.92\%$ on low sub-millimeter noise to $53.96\%$ on high noise, dampening high-frequency physiological flutter while maintaining responsiveness.

\subsection{Hysteresis State Stability Sweep (Benchmark B4)}
To evaluate gesture stability near the activation boundary, we simulated a stochastic state walk hovering within the decision boundary across 6 noise levels ($\sigma \in [0.015, 0.090]$) and 30 seeds ($180$ runs, $450\,000$ boundary frames total in \texttt{summary\_b4\_sweep.json}). Table~\ref{tab:hysteresis_sweep} compares a naive single-threshold detector ($\tau = 0.48$) against GestureCAD's dual-threshold automaton ($\tau_{\text{start}} = 0.44$, $\tau_{\text{release}} = 0.54$, debounced).

\begin{table}[t]
\caption{Gesture state chatter sensitivity sweep ($2\,500$ boundary frames per trial, $180$ total runs).}
\label{tab:hysteresis_sweep}
\centering
\small
\begin{tabular}{ccccc}
\toprule
\textbf{Noise $\sigma$} & \textbf{Naive Flips} & \textbf{Hysteresis Flips} & \textbf{Reduction} & \textbf{95\% CI} \\
\midrule
$0.015$ & $541.7 \pm 17.8$ & $3.1 \pm 0.4$ & $99.43\%$ & $[99.41\%, 99.46\%]$ \\
$0.030$ & $902.9 \pm 25.3$ & $26.2 \pm 5.6$ & $97.09\%$ & $[96.86\%, 97.32\%]$ \\
$0.045$ & $1067.1 \pm 26.3$ & $118.6 \pm 9.9$ & $88.88\%$ & $[88.53\%, 89.23\%]$ \\
$0.060$ & $1141.1 \pm 26.5$ & $217.1 \pm 16.7$ & $80.97\%$ & $[80.59\%, 81.35\%]$ \\
$0.075$ & $1173.3 \pm 24.5$ & $314.1 \pm 16.3$ & $73.23\%$ & $[72.84\%, 73.61\%]$ \\
$0.090$ & $1190.2 \pm 25.2$ & $405.0 \pm 15.3$ & $65.97\%$ & $[65.59\%, 66.35\%]$ \\
\midrule
\textbf{Grand Mean} & — & — & \textbf{85.55\%} & \textbf{[84.01\%, 87.10\%]} \\
\bottomrule
\end{tabular}
\end{table}

\textbf{Finding and Reconciliation}: 
\begin{itemize}
    \item In the comprehensive sweep, the dual-threshold automaton suppresses spurious state transitions by an aggregate mean of $85.55\%$ [$95\%$ CI: $84.01\%$–$87.10\%$, range: $65.6\%$–$99.4\%$]. Under typical resting tremor ($\sigma \le 0.030$), chatter elimination exceeds $97\%$.
    \item \textit{Reconciliation with Earlier Reports}: An isolated single-trial benchmark ($N=5\,000$ frames, seed 42) at $\sigma=0.045$ yielded a point-estimate reduction of $92.11\%$ (1,990 naive flips vs. 157 hysteresis flips in \texttt{summary\_b4\_hysteresis.json}). The 180-run parameter sweep reconciles this, showing that at $\sigma=0.045$, the multi-seed mean reduction is $88.88\%$ [$95\%$ CI: $88.53\%$–$89.23\%$], while the aggregate across all evaluated noise levels is $85.55\%$.
\end{itemize}

\subsection{Distance Scale Invariance (Benchmark B6)}
As derived analytically in Section~\ref{sec:methodology}, the normalized pinch metric $d_{\text{pinch}} = \|\mathbf{l}_4 - \mathbf{l}_8\|_2 / S_{\text{hand}}$ cancels camera depth $Z_{\text{cam}}$ under perspective projection. 

In analytical pinhole modeling across four distances ($30.0$, $50.0$, $80.0$, and $100.0\text{ cm}$) and three representative hand poses (Open Palm, Approaching Pinch, Contact Pinch in \texttt{summary\_b6\_scale\_invariance.json}), raw pixel distance varies by $70.0\%$ (coefficient of variation $\text{CV} = 70.0\%$), whereas the normalized metric maintains $\text{CV} = 0.0\%$. We reiterate that this $0.0\%$ invariance is an analytical mathematical property of perspective projection; physical camera trials with out-of-plane hand tilt introduce perspective foreshortening.

\subsection{Mesh Validity and Topological Invariants (Benchmark B9)}
We verified topological integrity across 20 extruded CAD models comprising regular prisms, cylinders, polygons ($N \in [3, 64]$), and non-convex profiles (L-shape, T-shape, U-shape, star 10-gon in \texttt{summary\_b9\_mesh\_validity.json}). Table~\ref{tab:mesh_validity} summarizes the results.

\begin{table}[t]
\caption{Topological invariant validation and volume verification across 20 geometric solid models.}
\label{tab:mesh_validity}
\centering
\small
\begin{tabular}{lcccc}
\toprule
\textbf{Model Category} & \textbf{Models} & \textbf{Euler ($\chi$)} & \textbf{Open Edges} & \textbf{Vol. Error} \\
\midrule
Rectangular Boxes & 5 & $2.000 \pm 0.0$ & 0 & $0.000\%$ \\
Cylinders ($N=32, 64$) & 5 & $2.000 \pm 0.0$ & 0 & $0.000\%$ \\
Regular Polygons (3 to 8-gon) & 5 & $2.000 \pm 0.0$ & 0 & $0.000\%$ \\
Non-Convex Profiles (L, T, U, Star) & 5 & $2.000 \pm 0.0$ & 0 & $0.000\%$ \\
\midrule
\textbf{Overall Compliance} & \textbf{20 / 20} & \textbf{100\% ($\chi=2$)} & \textbf{0 Boundary} & \textbf{$< 0.001\%$} \\
\bottomrule
\end{tabular}
\end{table}

\textbf{Finding}: All 20 models strictly satisfied:
\begin{enumerate}
    \item $\chi = V - E + F = 2.000$ ($100\%$ compliance).
    \item Zero unshared boundary edges ($\partial \Omega = \emptyset$, guaranteed watertightness).
    \item Zero non-manifold edges ($d(e) = 2$ for all edges).
    \item Exact Divergence Theorem volume calculation matching analytical formulas ($0.000\%$ numerical discrepancy).
\end{enumerate}
As established in Section~\ref{sec:geometry}, while non-convex models pass topological $\chi=2$ checks and Divergence Theorem volume integration algebraically, fan triangulation produces self-intersecting cap geometry on concave profiles. Non-convex geometries require Ear-Clipping (Earcut) to guarantee non-intersecting physical embeddings.

\subsection{Shape Regularization and Snapping Sweeps (B10, B11)}
\textbf{1. Heuristic Shape Regularization (B10)}: Evaluated across 6 distortion levels ($\sigma \in [0.04, 0.25]$) and 30 seeds ($18\,000$ synthetic strokes total in \texttt{summary\_b10\_sweep.json}). For low-to-moderate distortion ($\sigma \le 0.08$), primitive classification accuracy is $100.0\%$. At severe distortion ($\sigma = 0.20$ and $0.25$), accuracy degrades gracefully to $97.5\%$ and $92.1\%$, respectively.

\textbf{2. Geometric Snapping Precision (B11)}: Evaluated across 6 targeting dispersion radii ($r \in [0.10, 0.35]\text{ m}$) and 30 seeds ($36\,000$ queries total in \texttt{summary\_b11\_sweep.json}). Snapping captured an aggregate mean of $84.48\%$ [$95\%$ CI: $83.05\%$–$85.92\%$] of targeting attempts, reducing positional error by an aggregate mean of $79.74\%$ [$95\%$ CI: $77.28\%$–$82.20\%$, range: $41.7\%$–$95.7\%$]. (An earlier single-run benchmark with 500 queries at fixed low dispersion yielded $93.01\%$ error reduction in \texttt{summary\_b11\_snapping.json}; the 180-run sweep provides the robust multi-condition finding).

\begin{table*}[t]
\caption{Systematic component ablation study evaluating the quantitative impact of removing individual kinematic and geometric modules.}
\label{tab:ablation}
\centering
\small
\begin{tabular}{lccccc}
\toprule
\textbf{System Configuration} & \textbf{RMS Jitter} & \textbf{State Flip Red.} & \textbf{Targeting Err. Red.} & \textbf{Scale Invar. (CV)} & \textbf{Primary Failure Mode When Omitted} \\
\midrule
\textbf{Full GestureCAD System} & \textbf{1.251 mm} & \textbf{85.55\%} & \textbf{79.74\%} & \textbf{0.0\%} & \textit{Baseline operational performance across all modules.} \\
-- Hysteresis (Single Threshold $\tau=0.48$) & 1.251 mm & 0.0\% & 79.74\% & 0.0\% & Rapid limit-cycle pen chatter at decision boundary. \\
-- Adaptive Smoothing (Dead-Zone Only) & 3.176 mm & 85.55\% & 79.74\% & 0.0\% & Severe high-frequency hand tremor during active drawing strokes. \\
-- Dead-Zone (Static LERP Only) & 1.193 mm & 85.55\% & 79.74\% & 0.0\% & Micro-drift while holding steady; sluggish cursor response during sweeps. \\
-- Scale Normalization (Raw Pixels) & 1.251 mm & 85.55\% & 79.74\% & 70.0\% & Gestures fail when user moves between 30 cm and 100 cm camera depths. \\
-- Snapping (Raw Targeting) & 1.251 mm & 85.55\% & 0.0\% & 0.0\% & Inability to close sketch profiles cleanly; unclosed non-manifold boundaries. \\
\textbf{Unconditioned Baseline} & \textbf{3.176 mm} & \textbf{0.0\%} & \textbf{0.0\%} & \textbf{70.0\%} & Severe jitter, pen flutter, scale drift, and open sketch boundaries. \\
\bottomrule
\end{tabular}
\end{table*}

\subsection{Component Ablation Study}
Table~\ref{tab:ablation} synthesizes the quantitative contribution of each subsystem within GestureCAD. Disabling hysteresis eliminates chatter suppression ($0\%$ reduction), causing up to 1,990 spurious flips per 5,000 frames. Disabling adaptive smoothing exposes strokes to raw optical tremor ($3.176\text{ mm}$ RMS jitter). Omitting scale normalization introduces a $70\%$ variation in gesture thresholds across depths, and omitting snapping prevents users from reliably closing 2D sketch loops.

\section{Discussion}
\label{sec:discussion}

The development and empirical audit of GestureCAD provide several architectural, geometric, and human-factors insights into the integration of computer vision with geometric modeling.

\subsection{Architectural Decoupling and Web Performance}
A central design challenge in modern web-based graphics is managing the tension between declarative component state management (e.g., React, Vue) and high-frequency real-time rendering. Early iterations of web-based spatial interaction tools often routed tracking coordinates directly through component state variables. In our experiments, streaming 30--60 Hz landmark frames through React's Virtual DOM triggered continuous component re-renders, causing memory heap fragmentation and periodic garbage collection pauses ($\Delta t_{\text{gc}} > 25\text{ ms}$) that visibly disrupted drawing strokes.

GestureCAD's dual-path architecture resolves this tension:
\begin{enumerate}
    \item High-frequency spatial streams (cursor movements, live extrusion previews, plane manipulation) bypass the Virtual DOM entirely. Coordinates are maintained in mutable JavaScript references and written directly into Three.js \texttt{Float32Array} attribute buffers within the WebGL render loop (\texttt{useFrame}). This keeps cumulative software CAD execution under $0.2\text{ ms}$ per frame ($1.2\%$ of the $16.6\text{ ms}$ budget).
    \item Reactive state updates are reserved strictly for transactional boundaries (committing an extrusion, selecting a face, triggering Undo/Redo) managed via the Zustand store and Command Pattern.
\end{enumerate}
This architectural separation provides a practical blueprint for building high-performance spatial interaction systems within modern web frameworks.

\subsection{Topological Graph Validity vs. Spatial Geometric Embedding}
A key technical finding from our geometric audit is the vital distinction between \textit{topological graph invariance} and \textit{spatial geometric non-self-intersection}:
\begin{itemize}
    \item In Section~\ref{sec:geometry}, we showed that fan triangulation of planar caps yields a half-edge graph connectivity that strictly satisfies the Euler characteristic $\chi = V - E + F = 2$ and reports zero open boundary edges for any closed polygon.
    \item Furthermore, under Gauss's Divergence Theorem, overlapping exterior triangles contribute signed tetrahedral volumes of opposite sign that cancel algebraically, yielding exact agreement with analytical volume formulas even for non-convex profiles (as demonstrated in Benchmark B9).
    \item However, in physical $\mathbb{R}^3$ space, fan triangulation on concave profiles produces self-intersecting cap geometry. When exported to stereolithography (STL) files, self-intersecting facets violate the manifold assumptions of slicers and FEA meshers.
\end{itemize}
This insight underscores that topological validation ($\chi = 2$, $\partial \Omega = \emptyset$) is a necessary, but not sufficient, condition for geometric solid validity. Geometric modeling kernels must pair topological checks with explicit geometric non-self-intersection guarantees (e.g., via Ear-Clipping / Earcut triangulation).

\subsection{Practicality and Trade-Offs of Contactless CAD}
Contactless gesture CAD is not intended to replace high-precision desktop workstations for detailed engineering drafting. Rather, its practical utility emerges in specific contexts:
\begin{enumerate}
    \item \textbf{Sterile and Hazardous Environments}: In medical operating rooms (reviewing patient CT-scan geometry), pharmaceutical cleanrooms, or mechanical workshops with grease and debris, touching a mouse or keyboard contaminates the input peripheral. Contactless monocular webcam control provides an entirely hygienic interface.
    \item \textbf{Rapid Concept Ideation}: For industrial designers and architects, spatial hand gestures offer a fluid, natural medium for rough volumetric blocking and spatial proportion evaluation before committing to detailed drafting.
    \item \textbf{Ergonomic Diversity}: Sustained desktop mouse usage contributes to repetitive strain injuries (RSI) and carpal tunnel syndrome. GestureCAD provides an alternative physical modality for spatial manipulation.
\end{enumerate}
Conversely, the primary inherent trade-off of mid-air interaction is the \textbf{absence of haptic feedback}. When sketching on a physical tablet or using a mouse, physical surface friction stabilizes the human motor system. In mid-air, the lack of a physical backstop exacerbates physiological tremor. GestureCAD counterbalances this absence of tactile confirmation through software-assisted snapping, dead-zone gating, and clear visual state cues.

\subsection{Threats to Validity}
Several methodological factors must be considered when interpreting the reported results:
\begin{enumerate}
    \item \textbf{Synthetic vs. Real Noise Modeling}: While Benchmarks B3, B4, B10, and B11 evaluated comprehensive sweeps across 180 runs with multi-seed stochastic perturbations, the synthetic noise distributions (Gaussian perturbations and random-walk drift) may not fully capture the complex biomechanical dynamics of human motor tremor, fatigue-induced drift, or optical lighting flickers.
    \item \textbf{Hardware Dependency of Latency}: Software CAD latency (B1) was profiled on a high-performance host CPU (Intel Core i7-12700H). Performance may vary on lower-power client devices, mobile laptops, or under heavy browser background tab loads.
    \item \textbf{Unmeasured Motion-to-Photon Latency}: As discussed in Section~\ref{sec:experiments}, physical end-to-end latency remains unmeasured by high-speed optical hardware, representing an important gap between software execution time and perceptual responsiveness.
\end{enumerate}

\section{Limitations and Threats to Validity}
\label{sec:limitations}

To maintain rigorous scientific standards, this section explicitly outlines the current engineering and methodological boundaries of GestureCAD:

\begin{enumerate}
    \item \textbf{Absence of Formal Human-Subjects Evaluation}: While the system's kinematic filters, latency, and geometric validity have been verified empirically, formal human-subject trials evaluating task completion time, cognitive workload, and user satisfaction have not yet been executed. A standardized 12-participant within-subjects protocol (comparing GestureCAD against mouse-based CAD across 5 mechanical modeling tasks, accompanied by System Usability Scale [SUS]~\cite{brooke1996sus} and NASA-TLX~\cite{hart1988nasa} instruments) has been developed in the repository (\texttt{/study/}), but human recruitment and testing remain pending future work.
    \item \textbf{Unmeasured Physical Motion-to-Photon Latency}: Although the software CAD execution pipeline is verified at $0.1994\text{ ms}$ mean CPU time, total end-to-end physical motion-to-photon latency (incorporating optical exposure, USB bus transfer, neural inference, WebSocket IPC, WebGL rasterization, and monitor scanout) has not been measured with external high-speed optical hardware.
    \item \textbf{Topological Genus Constraint ($\chi = 2$)}: The current \texttt{SolidMesh} validation engine strictly enforces the genus-0 invariant $\chi = V - E + F = 2$. This topological check is valid for simply connected polyhedra homeomorphic to a 2-sphere (boxes, cylinders, convex prisms). Models containing through-holes, hollow cavities, or annular profiles possess genus $g \ge 1$ ($\chi \le 0$) and will be flagged as invalid. Generalizing the kernel requires supporting multi-loop 2D profiles with hole nesting and checking the generalized Euler-Poincar\'e formula $\chi = 2(1 - g)$.
    \item \textbf{Linear Extrusion Fan Cap Triangulation}: As analyzed in Section~\ref{sec:geometry}, \texttt{ExtrusionEngine} currently constructs polygon caps using a naive triangle fan from vertex 0. This method produces valid non-self-intersecting planar embeddings \textit{only for strictly convex profiles}. For concave profiles (e.g., L-shapes, star polygons), the resulting cap faces self-intersect in $\mathbb{R}^3$. Integrating Ear-Clipping (Earcut) triangulation---as already implemented in \texttt{RevolveEngine}---into \texttt{ExtrusionEngine} is necessary to support general simple polygons.
    \item \textbf{Monocular Depth Ambiguity and Hand Tilt}: While the wrist-to-MCP anatomical baseline provides distance scale invariance under frontal camera projection, severe out-of-plane hand tilt ($> 25^\circ$) introduces perspective foreshortening that affects 2D joint distance ratios. Incorporating temporal Kalman filtering or multi-view geometric constraints would enhance tilt robustness.
    \item \textbf{Absence of Constructive Solid Geometry (CSG) Booleans}: In the current prototype, intersecting geometric solids remain independent, overlapping polyhedral meshes; \texttt{BooleanEngine.js} is currently an unpopulated stub. Subtractive manufacturing operations (e.g., drilling holes, cutting pockets) require integrating a robust CSG boolean engine (e.g., via ManifoldCAD~\cite{manifoldcad} or OpenCASCADE.js~\cite{opencascadejs}).
    \item \textbf{Single-Hand Input Constraint}: The vision tracking pipeline is configured for a single hand (\texttt{max\_num\_hands=1}). Extending the architecture to bimanual input (e.g., non-dominant hand controlling camera orbit, pan, and zoom while dominant hand performs fine sketching) represents a natural ergonomic progression.
    \item \textbf{Client-Edge Architecture Dependency}: The system is not a 100% in-browser, zero-install application: it requires a running Python 3.10+ background daemon with OpenCV and MediaPipe. Migrating MediaPipe inference entirely to client-side WebAssembly (e.g., via TensorFlow.js or MediaPipe Web) will enable a true single-click, zero-install web application.
\end{enumerate}

\section{Conclusion and Future Work}
\label{sec:conclusion}

In this paper, we presented \textbf{GestureCAD}, a contactless, web-based 3D Computer-Aided Design modeling application driven by monocular RGB hand tracking. Rather than claiming an algorithmic breakthrough in computational geometry or machine learning, this work presents a practical systems-engineering integration that bridges commodity vision tracking with an interactive in-browser polyhedral geometry engine.

By establishing a decoupled client-edge architecture, GestureCAD separates Python-based neural hand tracking from a zero-reconciliation React Three Fiber WebGL client, executing software CAD operations within $0.1994\text{ ms}$ mean CPU time per frame on host hardware. To overcome the imprecision of contactless spatial input, we introduced a specialized kinematic conditioning pipeline: anatomical scale normalization eliminates camera depth variation under perspective projection, dual-threshold hysteresis suppresses boundary state chatter by an aggregate mean of $85.55\%$ [$95\%$ CI: $84.01\%$–$87.10\%$], velocity-adaptive smoothing attenuates RMS cursor jitter by $52.43\%$ [$95\%$ CI: $52.16\%$–$52.71\%$], and 3D snapping reduces targeting error by $79.74\%$ [$95\%$ CI: $77.28\%$–$82.20\%$].

Crucially, GestureCAD couples spatial interaction with explicit geometric validation. The system synthesizes closed polyhedral solids via linear extrusion and Rodrigues rotational sweeping, verifying 2-manifold topological invariants ($\chi = V - E + F = 2$, zero open boundary edges, zero non-manifold edges) and evaluating exact volumes via Gauss's Divergence Theorem. Furthermore, local $SO(3)$ orthonormal frame derivation enables hierarchical sketch-on-face operations directly on polyhedral boundaries and cylindrical tangent planes, with native export to Wavefront OBJ, AutoCAD DXF, and stereolithography STL.

\subsection{Future Work}
Building upon the verified engineering foundation established in this work, immediate avenues for future research include:
\begin{enumerate}
    \item \textbf{Arbitrary Polygon Triangulation}: Upgrading the linear extrusion engine from fan triangulation to robust Ear-Clipping (Earcut) to eliminate cap self-intersections on non-convex profiles.
    \item \textbf{Higher-Genus Topology Support}: Extending the geometric data structure to support profiles with nested interior holes, generalizing the topological invariant check to $\chi = 2(1 - g)$.
    \item \textbf{WebAssembly CSG Integration}: Incorporating a WebAssembly-compiled boolean evaluation engine (e.g., ManifoldCAD~\cite{manifoldcad} or OpenCASCADE.js~\cite{opencascadejs}) to support subtractive cutting, pocketing, and boolean unions.
    \item \textbf{Client-Side WASM Vision}: Migrating MediaPipe Hands inference entirely to client-side WebAssembly / WebGPU, transforming GestureCAD into a true zero-install, serverless web application.
    \item \textbf{Human-Factors Usability Evaluation}: Formally executing the 12-participant within-subjects usability study protocol developed in \texttt{/study/}, rigorously quantifying task completion times, error rates, SUS usability scores, and NASA-TLX cognitive workload distributions against standard mouse-driven desktop CAD.
\end{enumerate}
Through these developments, GestureCAD demonstrates a viable pathway toward making 3D spatial modeling more accessible, natural, and ergonomically adaptable without sacrificing the topological rigor demanded by modern engineering design.


\bibliographystyle{IEEEtran}
\bibliography{references}

\end{document}
